\documentclass[../../main.tex]{subfiles}

\begin{document}

\section{绝对连续函数}\label{section:ch06-绝对连续}

\begin{definition}[绝对连续]\label{definition:ch06-绝对连续}
闭有界区间 $[a,b]$ 上的实值函数 $f$ 称为在 $[a,b]$ 上是\textbf{绝对连续的}, 若对每个 $\varepsilon>0$, 存在 $\delta>0$, 使得对每个 $(a,b)$ 中的开区间的有限不交族 $\{(a_{k},b_{k})\}_{k=1}^{n}$,
\begin{align*}
\text{若 }\sum_{k=1}^{n}[b_{k}-a_{k}]<\delta,\quad\text{则 }\sum_{k=1}^{n}|f(b_{k})-f(a_{k})|<\varepsilon.
\end{align*}
\end{definition}

绝对连续性的准则在区间的有限族由单个区间组成的情形即为 $f$ 在 $[a,b]$ 上的一致连续性的准则. 因此绝对连续的函数是连续的. 反过来即使对递增函数也不成立.

\begin{example}\label{example:ch06-Cantor-Lebesgue}
Cantor-Lebesgue 函数 $\varphi$ 在 $[0,1]$ 上是递增与连续的, 但它不是绝对连续的 (也见习题 40 与 48). 事实上, 为看出 $\varphi$ 不是绝对连续的, 令 $n$ 为自然数. 在 Cantor 集构造的第 $n$ 步, $[0,1]$ 的 $2^{n}$ 个子区间的不交族 $\{[c_{k},d_{k}]\}_{1\leqslant k\leqslant 2^{n}}$ 被构造出, 它覆盖 Cantor 集, 每个区间的长度为 $(1/3)^{n}$. Cantor-Lebesgue 函数在该区间族由 $[0,1]$ 上的补集组成的每个区间上是常数. 因此, 由于 $\varphi$ 是递增的且 $\varphi(1)-\varphi(0)=1$,
\begin{align*}
\sum_{1\leqslant k\leqslant 2^{n}}[d_{k}-c_{k}]=(2/3)^{n}\quad\text{而}\quad\sum_{1\leqslant k\leqslant 2^{n}}[\varphi(d_{k})-\varphi(c_{k})]=1.
\end{align*}
不存在对应于关于 $\varphi$ 为绝对连续的准则的 $\varepsilon=1$ 的挑战.

显然绝对连续函数的线性组合是绝对连续的. 然而, 绝对连续函数的复合可以不是绝对连续的 (见习题 43、44 和 45).
\end{example}

\begin{proposition}\label{proposition:ch06-Lipschitz绝对连续}
若函数 $f$ 在闭有界区间 $[a,b]$ 上是 Lipschitz 的, 则它在 $[a,b]$ 上是绝对连续的.
\end{proposition}

\begin{proof}
令 $c>0$ 为 $f$ 在 $[a,b]$ 上的 Lipschitz 常数, 即
\begin{align*}
\text{对所有 }u,v\in[a,b],\quad|f(u)-f(v)|\leqslant c|u-v|.
\end{align*}
从而, 对于 $f$ 的绝对连续性的准则, 显然 $\delta=\varepsilon/c$ 对应于任何 $\varepsilon>0$ 的挑战.
\end{proof}

存在不是 Lipschitz 的绝对连续函数: $[0,1]$ 上定义为 $f(x)=\sqrt{x}$ 的函数 $f$ 是绝对连续的, 但不是 Lipschitz 的 (见习题 37).

\begin{theorem}\label{theorem:ch06-AC是递增差}
令函数 $f$ 在闭有界区间 $[a,b]$ 上是绝对连续的. 则 $f$ 是递增的绝对连续函数的差, 特别地, $f$ 是有界变差的.
\end{theorem}

\begin{proof}
我们首先证明 $f$ 是有界变差的. 事实上, 令 $\delta$ 对应于 $f$ 的绝对连续性的准则的 $\varepsilon=1$ 挑战. 令 $P$ 为将 $[a,b]$ 分为 $N$ 个闭区间, 每个区间的长度小于 $\delta$ 的分割 $\{c_{k},d_{k}\}_{k=1}^{N}$. 则根据关于 $f$ 的绝对连续性的 $\delta$ 的定义, 显然对于 $1\leqslant k\leqslant n$, $TV(f_{[c_{k},d_{k}]})\leqslant 1$. 加性公式\eqref{align:ch06-变差可加}推广到有限和. 因此,
\begin{align*}
TV(f)=\sum_{k=1}^{N}TV(f_{[c_{k},d_{k}]})\leqslant N.
\end{align*}
因此 $f$ 是有界变差的. 根据\eqref{align:ch06-直接表示}与绝对连续函数的和的绝对连续性, 为证明 $f$ 是递增的绝对连续函数的差, 仅须证明 $f$ 的全变差函数是绝对连续的. 令 $\varepsilon>0$. 选取 $\delta$ 对应于 $f$ 在 $[a,b]$ 上的绝对连续性的准则的 $\varepsilon/2$ 挑战. 令 $\{(c_{k},d_{k})\}_{k=1}^{n}$ 为 $(a,b)$ 的开子区间的不交族, 使得 $\sum_{k=1}^{n}[d_{k}-c_{k}]<\delta$. 对于 $1\leqslant k\leqslant n$, 令 $P_{k}$ 为 $[c_{k},d_{k}]$ 的一个分割. 根据 $f$ 在 $[a,b]$ 上的绝对连续性和 $\delta$ 的选取,
\begin{align*}
\sum_{k=1}^{n}V(f_{[c_{k},d_{k}]},P_{k})<\varepsilon/2.
\end{align*}
对 $1\leqslant k\leqslant n$, 取 $P_{k}$ 在 $[c_{k},d_{k}]$ 的分划变化的上确界, 得到
\begin{align*}
\sum_{k=1}^{n}TV(f_{[c_{k},d_{k}]})\leqslant\varepsilon/2<\varepsilon.
\end{align*}
我们从\eqref{align:ch06-变差递增}推出, 对 $1\leqslant k\leqslant n$, $TV(f_{[c_{k},d_{k}]})=TV(f_{[a,d_{k}]})-TV(f_{[a,c_{k}]})$. 因此
\begin{align}
\text{若 }\sum_{k=1}^{n}[d_{k}-c_{k}]<\delta,\text{则 }\sum_{k=1}^{n}|TV(f_{[a,d_{k}]})-TV(f_{[a,c_{k}]})|<\varepsilon.\label{align:ch06-变差AC}
\end{align}
因此 $f$ 的全变差函数在 $[a,b]$ 上是绝对连续的.
\end{proof}

\begin{theorem}\label{theorem:ch06-AC与一致可积}
令函数 $f$ 在闭有界区间 $[a,b]$ 上是连续的. 则 $f$ 在 $[a,b]$ 上是绝对连续的当且仅当均差函数族 $\{\operatorname{Diff}_{h}f\}_{0<h\leqslant 1}$ 在 $[a,b]$ 上是一致可积的.
\end{theorem}

\begin{proof}
首先假设 $\{\operatorname{Diff}_{h}f\}_{0<h\leqslant 1}$ 在 $[a,b]$ 上是一致可积的. 令 $\varepsilon>0$. 选取 $\delta>0$ 使得
\begin{align*}
\text{若 }m(E)<\delta\text{ 且 }0<h\leqslant 1,\quad\text{则}\int_{E}|\operatorname{Diff}_{h}f|<\varepsilon/2.
\end{align*}
我们宣称 $\delta$ 对应于关于 $f$ 为绝对连续的准则的 $\varepsilon$ 挑战. 事实上, 令 $\{(c_{k},d_{k})\}_{k=1}^{n}$ 为 $(a,b)$ 的不交子区间族使得 $\sum_{k=1}^{n}[d_{k}-c_{k}]<\delta$. 对于 $0<h\leqslant 1$ 与 $1\leqslant k\leqslant n$, 由\eqref{align:ch06-换元},
\begin{align*}
\operatorname{Avg}_{h}f(d_{k})-\operatorname{Avg}_{h}f(c_{k})=\int_{c_{k}}^{d_{k}}\operatorname{Diff}_{h}f.
\end{align*}
因此
\begin{align*}
\sum_{k=1}^{n}|\operatorname{Avg}_{h}f(d_{k})-\operatorname{Avg}_{h}f(c_{k})|\leqslant\sum_{k=1}^{n}\int_{c_{k}}^{d_{k}}|\operatorname{Diff}_{h}f|=\int_{E}|\operatorname{Diff}_{h}f|,
\end{align*}
其中 $E=\bigcup_{k=1}^{n}(c_{k},d_{k})$ 的测度小于 $\delta$. 因此, 根据 $\delta$ 的选取,
\begin{align*}
\text{对所有 }0<h\leqslant 1,\quad\sum_{k=1}^{n}|\operatorname{Avg}_{h}f(d_{k})-\operatorname{Avg}_{h}f(c_{k})|\leqslant\varepsilon/2.
\end{align*}
由于 $f$ 是连续的, 取 $h\to 0^{+}$ 的极限得到
\begin{align*}
\sum_{k=1}^{n}|f(d_{k})-f(c_{k})|\leqslant\varepsilon/2<\varepsilon.
\end{align*}
因此 $f$ 是绝对连续的.

为证明逆命题, 假定 $f$ 是绝对连续的. 前一个定理告诉我们 $f$ 是两个递增的绝对连续函数的差. 因此我们可以假设 $f$ 是递增的, 从而均差函数是非负的. 为证明 $\{\operatorname{Diff}_{h}\}_{0<h\leqslant 1}$ 的一致可积性, 令 $\varepsilon>0$. 我们必须证明存在 $\delta>0$, 使得对 $(a,b)$ 的每个可测子集 $E$, 若 $m(E)<\delta$ 且 $0<h\leqslant 1$, 则
\begin{align}
\int_{E}\operatorname{Diff}_{h}f<\varepsilon.\label{align:ch06-一致可积}
\end{align}
根据第 2 章的定理 11, 一个可测集 $E$ 包含于某个 $G_{\delta}$ 集 $G$, 其中 $m(G\sim E)=0$. 但每个 $G_{\delta}$ 集是开集的下降序列的交. 此外, 每个开集是开区间的可数族的不交并, 且因此每个开集是开集的上升序列的并, 该序列的每个集合是开区间的有限不交族的并. 因此, 根据积分的连续性, 为证明\eqref{align:ch06-一致可积}仅须找到 $\delta>0$ 使得对 $(a,b)$ 的开子区间的不交族 $\{(c_{k},d_{k})\}_{k=1}^{n}$, 若 $m(E)<\delta$, 则
\begin{align}
\int_{E}\operatorname{Diff}_{h}f<\varepsilon/2,\quad\text{其中 }E=\bigcup_{k=1}^{n}(c_{k},d_{k}),\quad\text{而 }0<h\leqslant 1.\label{align:ch06-有限和}
\end{align}
选取 $\delta>0$ 对应于 $f$ 在 $[a,b+1]$ 上的绝对连续性准则的 $\varepsilon/2$ 挑战. 通过 Riemann 积分的换元和抵消过程,
\begin{align*}
\int_{u}^{v}\operatorname{Diff}_{h}f=\frac{1}{h}\cdot\int_{0}^{h}g(t)\,dt,\quad\text{其中 }g(t)=f(v+t)-f(u+t),\quad 0\leqslant t\leqslant 1\text{ 而 }a\leqslant u<v\leqslant b.
\end{align*}
因此, 若 $\{(c_{k},d_{k})\}_{k=1}^{n}$ 是 $(a,b)$ 的开子区间的不交族,
\begin{align*}
\int_{E}\operatorname{Diff}_{h}f=\frac{1}{h}\cdot\int_{0}^{h}g(t)\,dt,
\end{align*}
其中
\begin{align*}
E=\bigcup_{k=1}^{n}(c_{k},d_{k}),\quad\text{而对所有 }0\leqslant t\leqslant 1,\quad g(t)=\sum_{k=1}^{n}[f(d_{k}+t)-f(c_{k}+t)].
\end{align*}
若 $\sum_{k=1}^{n}[d_{k}-c_{k}]<\delta$, 则对于 $0\leqslant t\leqslant 1$, $\sum_{k=1}^{n}[(d_{k}+t)-(c_{k}+t)]<\delta$, 且因此 $g(t)<\varepsilon/2$. 因此
\begin{align*}
\int_{E}\operatorname{Diff}_{h}f=\frac{1}{h}\cdot\int_{0}^{h}g(t)\,dt<\varepsilon/2.
\end{align*}
因而\eqref{align:ch06-有限和}对 $\delta$ 的这个选取成立.
\end{proof}

\begin{remark}\label{remark:ch06-函数族包含}
对非退化闭有界区间 $[a,b]$, 令 $\mathcal{F}_{Lip}$、$\mathcal{F}_{AC}$ 和 $\mathcal{F}_{BV}$ 分别表示 $[a,b]$ 上的 Lipschitz、绝对连续和有界变差函数族. 我们有以下严格包含关系:
\begin{align}
\mathcal{F}_{Lip}\subseteq\mathcal{F}_{AC}\subseteq\mathcal{F}_{BV}.\label{align:ch06-族包含}
\end{align}
命题\cref{proposition:ch06-Lipschitz绝对连续}告诉我们第一个包含关系, 而第二个包含关系在定理\cref{theorem:ch06-AC是递增差}中被证明. 这些族在线性组合的形式下封闭. 此外, 属于这些族中的某个族的函数, 在同一族中有其全变差函数, 因此, 根据\eqref{align:ch06-直接表示}, 属于这些族中的某个族的函数可表示为同族中的两个递增函数的差 (见习题 46).
\end{remark}

\begin{exercise}\label{exercise:ch06-37}
令 $f$ 为 $[0,1]$ 上的连续函数, 对每个 $0<\varepsilon<1$, 它在 $[\varepsilon,1]$ 上是绝对连续的.
\begin{enumerate}[(i)]
\item 证明 $f$ 在 $[0,1]$ 上可以不是绝对连续的.
\item 证明: 若 $f$ 是递增的, 则它在 $[0,1]$ 上是绝对连续的.
\item 证明 $[0,1]$ 上定义为 $f(x)=\sqrt{x}$ 的函数 $f$ 是绝对连续的, 但在 $[0,1]$ 上不是 Lipschitz 的.
\end{enumerate}
\end{exercise}

\begin{exercise}\label{exercise:ch06-38}
证明 $f$ 在 $[a,b]$ 上是绝对连续的当且仅当对每个 $\varepsilon>0$, 存在 $\delta>0$, 使得对 $(a,b)$ 内的每个开区间的可数不交族 $\{(a_{k},b_{k})\}_{k=1}^{\infty}$,
\begin{align*}
\text{若 }\sum_{k=1}^{n}[b_{k}-a_{k}]<\delta,\quad\text{则 }\sum_{k=1}^{n}|f(b_{k})-f(a_{k})|<\varepsilon.
\end{align*}
\end{exercise}

\begin{exercise}\label{exercise:ch06-39}
利用前一个习题证明: 若 $f$ 在 $[a,b]$ 上是递增的, 则 $f$ 在 $[a,b]$ 上是绝对连续的当且仅当对每个 $\varepsilon>0$, 存在 $\delta>0$ 使得对 $[a,b]$ 的可测子集 $E$,
\begin{align*}
\text{若 }m(E)<\delta,\quad m^{*}(f(E))<\varepsilon.
\end{align*}
\end{exercise}

\begin{exercise}\label{exercise:ch06-40}
用前一个习题证明 $[a,b]$ 上的递增的绝对连续函数 $f$ 将零测度集映上零测度集. 得出 Cantor-Lebesgue 函数 $\varphi$ 在 $[0,1]$ 上不是绝对连续的, 这是由于定义为 $\psi(x)=x+\varphi(x)$ ($0\leqslant x\leqslant 1$) 的函数 $\psi$ 将 Cantor 集映到测度为 $1$ 的集合 (见 2.7 节命题 21).
\end{exercise}

\begin{exercise}\label{exercise:ch06-41}
令 $f$ 为 $[a,b]$ 上递增的绝对连续函数. 用下面的 (i) 和 (ii) 得出 $f$ 将可测集映到可测集.
\begin{enumerate}[(i)]
\item 从 $f$ 的连续性和 $[a,b]$ 的紧性推出 $f$ 将闭集映到闭集, 因此将 $F_{\sigma}$ 集映到 $F_{\sigma}$ 集.
\item 前一个习题告诉我们 $f$ 将零测度集映上零测度集.
\end{enumerate}
\end{exercise}

\begin{exercise}\label{exercise:ch06-42}
证明绝对连续函数的和与积都是绝对连续的.
\end{exercise}

\begin{exercise}\label{exercise:ch06-43}
定义 $[-1,1]$ 上的函数 $f$ 和 $g$ 为 $f(x)=x^{\frac{1}{3}}$ ($-1\leqslant x\leqslant 1$) 与
\begin{align*}
g(x)=
\begin{cases}
x^{2}\cos(\pi/2x)&\text{若 }x\neq 0,x\in[-1,1]\\
0&\text{若 }x=0
\end{cases}
\end{align*}
\begin{enumerate}[(i)]
\item 证明 $f$ 和 $g$ 在 $[-1,1]$ 上都是绝对连续的.
\item 对于 $[-1,1]$ 的分割 $P_{n}=\{-1,0,1/2n,1/[2n-1],\cdots,1/3,1/2,1\}$, 讨论 $V(f\circ g,P_{n})$.
\item 证明 $f\circ g$ 不是有界变差的, 且因此在 $[-1,1]$ 上也不是绝对连续的.
\end{enumerate}
\end{exercise}

\begin{exercise}\label{exercise:ch06-44}
令 $f$ 在 $\mathbf{R}$ 上是 Lipschitz 的而 $g$ 在 $[a,b]$ 上是绝对连续的. 证明复合 $f\circ g$ 在 $[a,b]$ 上是绝对连续的.
\end{exercise}

\begin{exercise}\label{exercise:ch06-45}
令 $f$ 在 $\mathbf{R}$ 上是绝对连续的而 $g$ 在 $[a,b]$ 上是绝对连续与严格单调的. 证明复合 $f\circ g$ 在 $[a,b]$ 上是绝对连续的.
\end{exercise}

\begin{exercise}\label{exercise:ch06-46}
验证本节最后注的断言.
\end{exercise}

\begin{exercise}\label{exercise:ch06-47}
证明 $f$ 在 $[a,b]$ 上是绝对连续的当且仅当对每个 $\varepsilon>0$, 存在 $\delta>0$, 使得对 $(a,b)$ 中的开区间的有限不交族 $\{(a_{k},b_{k})\}_{k=1}^{n}$,
\begin{align*}
\text{若 }\sum_{k=1}^{n}[b_{k}-a_{k}]<\delta,\quad\left|\sum_{k=1}^{n}[f(b_{k})-f(a_{k})]\right|<\varepsilon.
\end{align*}
\end{exercise}

\end{document}